
\documentclass{article}
\usepackage{colortbl}
\usepackage{makecell}
\usepackage{multirow}
\usepackage{supertabular}

\begin{document}

\newcounter{utterance}

\centering \large Interaction Transcript for game `clean\_up', experiment `0\_easy\_3obj\_fr', episode 1 with Apertus{-}v1.5{-}8B.
\vspace{24pt}

{ \footnotesize  \setcounter{utterance}{1}
\setlength{\tabcolsep}{0pt}
\begin{supertabular}{c@{$\;$}|p{.15\linewidth}@{}p{.15\linewidth}p{.15\linewidth}p{.15\linewidth}p{.15\linewidth}p{.15\linewidth}}
    \# & $\;$A & \multicolumn{4}{c}{Game Master} & $\;\:$B\\
    \hline

    \theutterance \stepcounter{utterance}  
    & & \multicolumn{4}{p{0.6\linewidth}}{
        \cellcolor[rgb]{0.9,0.9,0.9}{
            \makecell[{{p{\linewidth}}}]{
                \texttt{\tiny{[P1$\langle$GM]}}
                \texttt{Je suis votre maître de jeu, et vous jouez à un jeu collaboratif avec la grille suivante comme plateau de jeu :} \\
\\ 
\texttt{```} \\
\texttt{ 1234567} \\
\texttt{╔═══════╗} \\
\texttt{║◌◌◌◌◌◌◌║ 1} \\
\texttt{║◌C◌◌◌◌◌║ 2} \\
\texttt{║◌◌◌◌◌┌─╢ 3} \\
\texttt{║◌◌◌┌─┤◌║ 4} \\
\texttt{║◌◌◌│L├─╢ 5} \\
\texttt{╟───┼─┘◌║ 6} \\
\texttt{║◌◌◌│◌◌P║ 7} \\
\texttt{╚═══╧═══╝} \\
\\ 
\texttt{```} \\
\\ 
\texttt{* Le bord supérieur affiche des coordonnées x croissantes vers la droite, et le bord droit des coordonnées y croissantes vers le bas.} \\
\texttt{* Les objets suivants sont placés aléatoirement sur votre grille : 'C', 'L', 'P'.} \\
\\ 
\\ 
\texttt{L'autre joueur voit une variante du plateau de jeu, où les objets sont placés à des emplacements aléatoires différents. Vous ne pouvez pas voir le plateau de l'autre joueur et il ne peut pas voir le vôtre.} \\
\\ 
\\ 
\texttt{**Objectif :**} \\
\\ 
\texttt{Les deux joueurs doivent déplacer les objets sur leur arrière{-}plan respectif de manière à ce que des objets identiques se retrouvent aux mêmes coordonnées. Vous devez communiquer avec l'autre joueur pour convenir d'une configuration d'objectif commune.} \\
\\ 
\\ 
\texttt{**Règles :**} \\
\\ 
\texttt{* À chaque tour, vous pouvez envoyer exactement l'une des deux commandes suivantes :} \\
\texttt{ 1. `DIRE: <MESSAGE>` : pour envoyer un message (tout ce qui se trouve jusqu'au prochain saut de ligne) à l'autre joueur. Je le transmettrai à votre partenaire.} \\
\texttt{2. `DÉPLACER: <OBJET>, (<X>, <Y>)` : pour déplacer un objet vers une nouvelle position, où `<X>` est la colonne et `<Y>` est la ligne. Je vous informerai si votre déplacement a réussi ou non.} \\
\texttt{* Si vous ne respectez pas le format, ou si vous envoyez plusieurs commandes à la fois, je dois vous pénaliser.} \\
\texttt{* Si les deux joueurs accumulent plus de 8 pénalités, vous perdez tous les deux la partie.} \\
\texttt{* Il est essentiel que vous communiquiez avec l'autre joueur concernant votre état d'objectif ! La *seule* façon de transmettre votre stratégie à l'autre joueur est d'utiliser la commande `DIRE: <MESSAGE>` !} \\
\\ 
\\ 
\texttt{**Déplacement d'objets**} \\
\\ 
\texttt{* Vous ne pouvez déplacer des objets que vers des cellules situées à l'intérieur des limites de la grille. La cellule cible doit être vide, c'est{-}à{-}dire qu'elle ne doit contenir que le symbole '◌'.} \\
\texttt{* Si vous essayez de déplacer un objet vers un endroit qui n'est pas vide, ou si vous essayez de le déplacer en dehors de la grille, je dois vous pénaliser. Vous avez droit à un nouvel essai.} \\
\texttt{* Avant d'effectuer un déplacement, vérifiez bien que l'emplacement cible est vide et qu'il ne contient aucune lettre, aucun cadre ou aucune ligne !} \\
\\ 
\\ 
\texttt{**Fin de partie**} \\
\\ 
\texttt{Si vous pensez avoir atteint l'objectif d'aligner tous les objets, vous pouvez demander à l'autre joueur de terminer la partie en envoyant `DIRE: terminé ?`. Si l'autre joueur vous demande de terminer la partie et que vous répondez `DIRE: terminé !`, la partie se termine.} \\
\\ 
\texttt{Les deux joueurs gagnent si la partie se termine en 12 tours, un tour étant défini comme l'envoi d'une commande valide par chacun des deux joueurs.} \\
\\ 
\texttt{**Score :**} \\
\\ 
\texttt{Plus les objets identiques sont proches les uns des autres sur les deux plateaux de jeu, plus vous obtenez de points. Les pénalités réduisent vos points. Pouvez{-}vous battre le record ?} \\
\\ 
\\ 
\texttt{Envoyez un message à l'autre joueur pour commencer le jeu !} \\
            }
        }
    }
    & \\ \\

    \theutterance \stepcounter{utterance}  
    & \multicolumn{4}{p{0.6\linewidth}}{
        \cellcolor[rgb]{0.9,0.9,0.9}{
            \makecell[{{p{\linewidth}}}]{
                \texttt{\tiny{[P1$\rangle$GM]}}
                \texttt{DIRE: Bonjour, commençons le jeu. Je vois que nous devons aligner les objets 'C', 'L' et 'P'. Est{-}ce que vous avez une idée de où pourrait être placé l'autre objet ? Je vais commencer à déplacer le 'C' vers la position (2, 3).} \\
            }
        }
    }
    & & \\ \\

    \theutterance \stepcounter{utterance}  
    & & & \multicolumn{2}{p{0.3\linewidth}}{
        \cellcolor[rgb]{0.9,0.9,0.9}{
            \makecell[{{p{\linewidth}}}]{
                \texttt{\tiny{[GM$|$GM]}}
                \texttt{DIRE: Bonjour, commençons le jeu. Je vois que nous devons aligner les objets 'C', 'L' et 'P'. Est{-}ce que vous avez une idée de où pourrait être placé l'autre objet ? Je vais commencer à déplacer le 'C' vers la position (2, 3).} \\
            }
        }
    }
    & & \\ \\

    \theutterance \stepcounter{utterance}  
    & & \multicolumn{4}{p{0.6\linewidth}}{
        \cellcolor[rgb]{0.9,0.9,0.9}{
            \makecell[{{p{\linewidth}}}]{
                \texttt{\tiny{[GM$\rangle$P2]}}
                \texttt{Je suis votre maître de jeu, et vous jouez à un jeu collaboratif avec la grille suivante comme plateau de jeu :} \\
\\ 
\texttt{```} \\
\texttt{ 1234567} \\
\texttt{╔═══════╗} \\
\texttt{║◌◌◌◌◌◌◌║ 1} \\
\texttt{║◌◌C◌◌◌◌║ 2} \\
\texttt{║◌◌◌◌◌┌─╢ 3} \\
\texttt{║◌◌◌┌─┤◌║ 4} \\
\texttt{║◌P◌│◌├─╢ 5} \\
\texttt{╟───┼─┘◌║ 6} \\
\texttt{║L◌◌│◌◌◌║ 7} \\
\texttt{╚═══╧═══╝} \\
\\ 
\texttt{```} \\
\\ 
\texttt{* Le bord supérieur affiche des coordonnées x croissantes vers la droite, et le bord droit des coordonnées y croissantes vers le bas.} \\
\texttt{* Les objets suivants sont placés aléatoirement sur votre grille : 'C', 'L', 'P'.} \\
\\ 
\\ 
\texttt{L'autre joueur voit une variante du plateau de jeu, où les objets sont placés à des emplacements aléatoires différents. Vous ne pouvez pas voir le plateau de l'autre joueur et il ne peut pas voir le vôtre.} \\
\\ 
\\ 
\texttt{**Objectif :**} \\
\\ 
\texttt{Les deux joueurs doivent déplacer les objets sur leur arrière{-}plan respectif de manière à ce que des objets identiques se retrouvent aux mêmes coordonnées. Vous devez communiquer avec l'autre joueur pour convenir d'une configuration d'objectif commune.} \\
\\ 
\\ 
\texttt{**Règles :**} \\
\\ 
\texttt{* À chaque tour, vous pouvez envoyer exactement l'une des deux commandes suivantes :} \\
\texttt{ 1. `DIRE: <MESSAGE>` : pour envoyer un message (tout ce qui se trouve jusqu'au prochain saut de ligne) à l'autre joueur. Je le transmettrai à votre partenaire.} \\
\texttt{2. `DÉPLACER: <OBJET>, (<X>, <Y>)` : pour déplacer un objet vers une nouvelle position, où `<X>` est la colonne et `<Y>` est la ligne. Je vous informerai si votre déplacement a réussi ou non.} \\
\texttt{* Si vous ne respectez pas le format, ou si vous envoyez plusieurs commandes à la fois, je dois vous pénaliser.} \\
\texttt{* Si les deux joueurs accumulent plus de 8 pénalités, vous perdez tous les deux la partie.} \\
\texttt{* Il est essentiel que vous communiquiez avec l'autre joueur concernant votre état d'objectif ! La *seule* façon de transmettre votre stratégie à l'autre joueur est d'utiliser la commande `DIRE: <MESSAGE>` !} \\
\\ 
\\ 
\texttt{**Déplacement d'objets**} \\
\\ 
\texttt{* Vous ne pouvez déplacer des objets que vers des cellules situées à l'intérieur des limites de la grille. La cellule cible doit être vide, c'est{-}à{-}dire qu'elle ne doit contenir que le symbole '◌'.} \\
\texttt{* Si vous essayez de déplacer un objet vers un endroit qui n'est pas vide, ou si vous essayez de le déplacer en dehors de la grille, je dois vous pénaliser. Vous avez droit à un nouvel essai.} \\
\texttt{* Avant d'effectuer un déplacement, vérifiez bien que l'emplacement cible est vide et qu'il ne contient aucune lettre, aucun cadre ou aucune ligne !} \\
\\ 
\\ 
\texttt{**Fin de partie**} \\
\\ 
\texttt{Si vous pensez avoir atteint l'objectif d'aligner tous les objets, vous pouvez demander à l'autre joueur de terminer la partie en envoyant `DIRE: terminé ?`. Si l'autre joueur vous demande de terminer la partie et que vous répondez `DIRE: terminé !`, la partie se termine.} \\
\\ 
\texttt{Les deux joueurs gagnent si la partie se termine en 12 tours, un tour étant défini comme l'envoi d'une commande valide par chacun des deux joueurs.} \\
\\ 
\texttt{**Score :**} \\
\\ 
\texttt{Plus les objets identiques sont proches les uns des autres sur les deux plateaux de jeu, plus vous obtenez de points. Les pénalités réduisent vos points. Pouvez{-}vous battre le record ?} \\
\\ 
\texttt{L'autre joueur a commencé le jeu en envoyant ce message :} \\
\\ 
\texttt{"Bonjour, commençons le jeu. Je vois que nous devons aligner les objets 'C', 'L' et 'P'. Est{-}ce que vous avez une idée de où pourrait être placé l'autre objet ? Je vais commencer à déplacer le 'C' vers la position (2, 3)."} \\
\\ 
\texttt{Quelle est votre première commande ?} \\
            }
        }
    }
    & \\ \\

    \theutterance \stepcounter{utterance}  
    & & & \multicolumn{4}{p{0.6\linewidth}}{
        \cellcolor[rgb]{0.9,0.9,0.9}{
            \makecell[{{p{\linewidth}}}]{
                \texttt{\tiny{[GM$\langle$P2]}}
                \texttt{DÉPLACER: C, (2, 3)} \\
            }
        }
    }
     \\ \\

    \theutterance \stepcounter{utterance}  
    & & & \multicolumn{2}{p{0.3\linewidth}}{
        \cellcolor[rgb]{0.9,0.9,0.9}{
            \makecell[{{p{\linewidth}}}]{
                \texttt{\tiny{[GM$|$GM]}}
                \texttt{DÉPLACER: C, (2, 3)} \\
            }
        }
    }
    & & \\ \\

    \theutterance \stepcounter{utterance}  
    & & & \multicolumn{2}{p{0.3\linewidth}}{
        \cellcolor[rgb]{0.9,0.9,0.9}{
            \makecell[{{p{\linewidth}}}]{
                \texttt{\tiny{[GM$|$GM]}}
                \texttt{C moved to (2, 3)} \\
            }
        }
    }
    & & \\ \\

    \theutterance \stepcounter{utterance}  
    & & \multicolumn{4}{p{0.6\linewidth}}{
        \cellcolor[rgb]{0.9,0.9,0.9}{
            \makecell[{{p{\linewidth}}}]{
                \texttt{\tiny{[P1$\langle$GM]}}
                \texttt{Votre message a été transmis à l'autre joueur.} \\
\\ 
\\ 
\texttt{Vous jouez actuellement le tour 1 sur un maximum de 12.} \\
\texttt{Vous avez collectivement accumulé 0 pénalités sur 8.} \\
\texttt{L'autre joueur a déplacé un objet sur sa grille.} \\
\\ 
\\ 
\texttt{Quelle est votre prochaine commande ?} \\
            }
        }
    }
    & \\ \\

    \theutterance \stepcounter{utterance}  
    & \multicolumn{4}{p{0.6\linewidth}}{
        \cellcolor[rgb]{0.9,0.9,0.9}{
            \makecell[{{p{\linewidth}}}]{
                \texttt{\tiny{[P1$\rangle$GM]}}
                \texttt{DIRE: Avez{-}vous pu déplacer l'un des objets sur votre grille ? Si oui, pouvez{-}vous me dire où se trouve votre objet 'P' ? Je vais déplacer mon 'C' vers (3, 2) si c'est approprié.} \\
            }
        }
    }
    & & \\ \\

    \theutterance \stepcounter{utterance}  
    & & & \multicolumn{2}{p{0.3\linewidth}}{
        \cellcolor[rgb]{0.9,0.9,0.9}{
            \makecell[{{p{\linewidth}}}]{
                \texttt{\tiny{[GM$|$GM]}}
                \texttt{DIRE: Avez{-}vous pu déplacer l'un des objets sur votre grille ? Si oui, pouvez{-}vous me dire où se trouve votre objet 'P' ? Je vais déplacer mon 'C' vers (3, 2) si c'est approprié.} \\
            }
        }
    }
    & & \\ \\

    \theutterance \stepcounter{utterance}  
    & & \multicolumn{4}{p{0.6\linewidth}}{
        \cellcolor[rgb]{0.9,0.9,0.9}{
            \makecell[{{p{\linewidth}}}]{
                \texttt{\tiny{[GM$\rangle$P2]}}
                \texttt{Déplacement de 'C' vers (2,3) réussi. Votre grille mise à jour ressemble à ceci :} \\
\texttt{```} \\
\texttt{ 1234567} \\
\texttt{╔═══════╗} \\
\texttt{║◌◌◌◌◌◌◌║ 1} \\
\texttt{║◌◌◌◌◌◌◌║ 2} \\
\texttt{║◌C◌◌◌┌─╢ 3} \\
\texttt{║◌◌◌┌─┤◌║ 4} \\
\texttt{║◌P◌│◌├─╢ 5} \\
\texttt{╟───┼─┘◌║ 6} \\
\texttt{║L◌◌│◌◌◌║ 7} \\
\texttt{╚═══╧═══╝} \\
\texttt{```} \\
\texttt{Vous jouez actuellement le tour 2 sur un maximum de 12.} \\
\texttt{Vous avez collectivement accumulé 0 pénalités sur 8.} \\
\texttt{L'autre joueur a envoyé ce message :} \\
\\ 
\\ 
\texttt{"Avez{-}vous pu déplacer l'un des objets sur votre grille ? Si oui, pouvez{-}vous me dire où se trouve votre objet 'P' ? Je vais déplacer mon 'C' vers (3, 2) si c'est approprié."} \\
\\ 
\\ 
\texttt{Quelle est votre prochaine commande ?} \\
            }
        }
    }
    & \\ \\

    \theutterance \stepcounter{utterance}  
    & & & \multicolumn{4}{p{0.6\linewidth}}{
        \cellcolor[rgb]{0.9,0.9,0.9}{
            \makecell[{{p{\linewidth}}}]{
                \texttt{\tiny{[GM$\langle$P2]}}
                \texttt{DIRE: Mon objet 'P' se trouve actuellement à la position (4, 5).} \\
            }
        }
    }
     \\ \\

    \theutterance \stepcounter{utterance}  
    & & & \multicolumn{2}{p{0.3\linewidth}}{
        \cellcolor[rgb]{0.9,0.9,0.9}{
            \makecell[{{p{\linewidth}}}]{
                \texttt{\tiny{[GM$|$GM]}}
                \texttt{DIRE: Mon objet 'P' se trouve actuellement à la position (4, 5).} \\
            }
        }
    }
    & & \\ \\

    \theutterance \stepcounter{utterance}  
    & & \multicolumn{4}{p{0.6\linewidth}}{
        \cellcolor[rgb]{0.9,0.9,0.9}{
            \makecell[{{p{\linewidth}}}]{
                \texttt{\tiny{[P1$\langle$GM]}}
                \texttt{Votre message a été transmis à l'autre joueur.} \\
\\ 
\\ 
\texttt{Vous jouez actuellement le tour 2 sur un maximum de 12.} \\
\texttt{Vous avez collectivement accumulé 0 pénalités sur 8.} \\
\texttt{L'autre joueur a envoyé ce message :} \\
\\ 
\\ 
\texttt{"Mon objet 'P' se trouve actuellement à la position (4, 5)."} \\
\\ 
\\ 
\texttt{Quelle est votre prochaine commande ?} \\
            }
        }
    }
    & \\ \\

    \theutterance \stepcounter{utterance}  
    & \multicolumn{4}{p{0.6\linewidth}}{
        \cellcolor[rgb]{0.9,0.9,0.9}{
            \makecell[{{p{\linewidth}}}]{
                \texttt{\tiny{[P1$\rangle$GM]}}
                \texttt{DIRE: Merci pour l'information. Je vais déplacer mon 'C' vers (2, 3) pour être plus proche du 'C' potentiel de l'autre joueur. Veuillez vérifier si cette position est libre sur votre grille.} \\
            }
        }
    }
    & & \\ \\

    \theutterance \stepcounter{utterance}  
    & & & \multicolumn{2}{p{0.3\linewidth}}{
        \cellcolor[rgb]{0.9,0.9,0.9}{
            \makecell[{{p{\linewidth}}}]{
                \texttt{\tiny{[GM$|$GM]}}
                \texttt{DIRE: Merci pour l'information. Je vais déplacer mon 'C' vers (2, 3) pour être plus proche du 'C' potentiel de l'autre joueur. Veuillez vérifier si cette position est libre sur votre grille.} \\
            }
        }
    }
    & & \\ \\

    \theutterance \stepcounter{utterance}  
    & & \multicolumn{4}{p{0.6\linewidth}}{
        \cellcolor[rgb]{0.9,0.9,0.9}{
            \makecell[{{p{\linewidth}}}]{
                \texttt{\tiny{[GM$\rangle$P2]}}
                \texttt{Votre message a été transmis à l'autre joueur.} \\
\\ 
\\ 
\texttt{Vous jouez actuellement le tour 3 sur un maximum de 12.} \\
\texttt{Vous avez collectivement accumulé 0 pénalités sur 8.} \\
\texttt{L'autre joueur a envoyé ce message :} \\
\\ 
\\ 
\texttt{"Merci pour l'information. Je vais déplacer mon 'C' vers (2, 3) pour être plus proche du 'C' potentiel de l'autre joueur. Veuillez vérifier si cette position est libre sur votre grille."} \\
\\ 
\\ 
\texttt{Quelle est votre prochaine commande ?} \\
            }
        }
    }
    & \\ \\

    \theutterance \stepcounter{utterance}  
    & & & \multicolumn{4}{p{0.6\linewidth}}{
        \cellcolor[rgb]{0.9,0.9,0.9}{
            \makecell[{{p{\linewidth}}}]{
                \texttt{\tiny{[GM$\langle$P2]}}
                \texttt{DÉPLACER: C, (2, 3)} \\
            }
        }
    }
     \\ \\

    \theutterance \stepcounter{utterance}  
    & & & \multicolumn{2}{p{0.3\linewidth}}{
        \cellcolor[rgb]{0.9,0.9,0.9}{
            \makecell[{{p{\linewidth}}}]{
                \texttt{\tiny{[GM$|$GM]}}
                \texttt{DÉPLACER: C, (2, 3)} \\
            }
        }
    }
    & & \\ \\

    \theutterance \stepcounter{utterance}  
    & & & \multicolumn{2}{p{0.3\linewidth}}{
        \cellcolor[rgb]{0.9,0.9,0.9}{
            \makecell[{{p{\linewidth}}}]{
                \texttt{\tiny{[GM$|$GM]}}
                \texttt{Pénalité : (2,3) n'est pas vide, mais contient 'C'.} \\
\\ 
\texttt{Vous avez collectivement accumulé 1 pénalités sur 8.} \\
\texttt{Veuillez réessayer !} \\
            }
        }
    }
    & & \\ \\

    \theutterance \stepcounter{utterance}  
    & & \multicolumn{4}{p{0.6\linewidth}}{
        \cellcolor[rgb]{0.9,0.9,0.9}{
            \makecell[{{p{\linewidth}}}]{
                \texttt{\tiny{[GM$\rangle$P2]}}
                \texttt{Pénalité : (2,3) n'est pas vide, mais contient 'C'.} \\
\\ 
\texttt{Vous avez collectivement accumulé 1 pénalités sur 8.} \\
\texttt{Veuillez réessayer !} \\
            }
        }
    }
    & \\ \\

    \theutterance \stepcounter{utterance}  
    & & & \multicolumn{4}{p{0.6\linewidth}}{
        \cellcolor[rgb]{0.9,0.9,0.9}{
            \makecell[{{p{\linewidth}}}]{
                \texttt{\tiny{[GM$\langle$P2]}}
                \texttt{DIRE: Je vais vérifier la position vacante pour mon 'C'. Puis{-}je demander à l'autre joueur de me donner les coordonnées disponibles pour que je puisse effectuer un déplacement valide ?} \\
            }
        }
    }
     \\ \\

    \theutterance \stepcounter{utterance}  
    & & & \multicolumn{2}{p{0.3\linewidth}}{
        \cellcolor[rgb]{0.9,0.9,0.9}{
            \makecell[{{p{\linewidth}}}]{
                \texttt{\tiny{[GM$|$GM]}}
                \texttt{DIRE: Je vais vérifier la position vacante pour mon 'C'. Puis{-}je demander à l'autre joueur de me donner les coordonnées disponibles pour que je puisse effectuer un déplacement valide ?} \\
            }
        }
    }
    & & \\ \\

    \theutterance \stepcounter{utterance}  
    & & \multicolumn{4}{p{0.6\linewidth}}{
        \cellcolor[rgb]{0.9,0.9,0.9}{
            \makecell[{{p{\linewidth}}}]{
                \texttt{\tiny{[P1$\langle$GM]}}
                \texttt{Votre message a été transmis à l'autre joueur.} \\
\\ 
\\ 
\texttt{Vous jouez actuellement le tour 3 sur un maximum de 12.} \\
\texttt{Vous avez collectivement accumulé 1 pénalités sur 8.} \\
\texttt{L'autre joueur a envoyé ce message :} \\
\\ 
\\ 
\texttt{"Je vais vérifier la position vacante pour mon 'C'. Puis{-}je demander à l'autre joueur de me donner les coordonnées disponibles pour que je puisse effectuer un déplacement valide ?"} \\
\\ 
\\ 
\texttt{Quelle est votre prochaine commande ?} \\
            }
        }
    }
    & \\ \\

    \theutterance \stepcounter{utterance}  
    & \multicolumn{4}{p{0.6\linewidth}}{
        \cellcolor[rgb]{0.9,0.9,0.9}{
            \makecell[{{p{\linewidth}}}]{
                \texttt{\tiny{[P1$\rangle$GM]}}
                \texttt{DIRE: Bien sûr, je peux vous donner les coordonnées disponibles. Sur ma grille, la position (2, 3) est actuellement vide. Veuillez déplacer votre 'C' vers cette position si c'est possible selon vos règles.} \\
            }
        }
    }
    & & \\ \\

    \theutterance \stepcounter{utterance}  
    & & & \multicolumn{2}{p{0.3\linewidth}}{
        \cellcolor[rgb]{0.9,0.9,0.9}{
            \makecell[{{p{\linewidth}}}]{
                \texttt{\tiny{[GM$|$GM]}}
                \texttt{DIRE: Bien sûr, je peux vous donner les coordonnées disponibles. Sur ma grille, la position (2, 3) est actuellement vide. Veuillez déplacer votre 'C' vers cette position si c'est possible selon vos règles.} \\
            }
        }
    }
    & & \\ \\

    \theutterance \stepcounter{utterance}  
    & & \multicolumn{4}{p{0.6\linewidth}}{
        \cellcolor[rgb]{0.9,0.9,0.9}{
            \makecell[{{p{\linewidth}}}]{
                \texttt{\tiny{[GM$\rangle$P2]}}
                \texttt{Votre message a été transmis à l'autre joueur.} \\
\\ 
\\ 
\texttt{Vous jouez actuellement le tour 4 sur un maximum de 12.} \\
\texttt{Vous avez collectivement accumulé 1 pénalités sur 8.} \\
\texttt{L'autre joueur a envoyé ce message :} \\
\\ 
\\ 
\texttt{"Bien sûr, je peux vous donner les coordonnées disponibles. Sur ma grille, la position (2, 3) est actuellement vide. Veuillez déplacer votre 'C' vers cette position si c'est possible selon vos règles."} \\
\\ 
\\ 
\texttt{Quelle est votre prochaine commande ?} \\
            }
        }
    }
    & \\ \\

    \theutterance \stepcounter{utterance}  
    & & & \multicolumn{4}{p{0.6\linewidth}}{
        \cellcolor[rgb]{0.9,0.9,0.9}{
            \makecell[{{p{\linewidth}}}]{
                \texttt{\tiny{[GM$\langle$P2]}}
                \texttt{DÉPLACER: C, (2, 3)} \\
            }
        }
    }
     \\ \\

    \theutterance \stepcounter{utterance}  
    & & & \multicolumn{2}{p{0.3\linewidth}}{
        \cellcolor[rgb]{0.9,0.9,0.9}{
            \makecell[{{p{\linewidth}}}]{
                \texttt{\tiny{[GM$|$GM]}}
                \texttt{DÉPLACER: C, (2, 3)} \\
            }
        }
    }
    & & \\ \\

    \theutterance \stepcounter{utterance}  
    & & & \multicolumn{2}{p{0.3\linewidth}}{
        \cellcolor[rgb]{0.9,0.9,0.9}{
            \makecell[{{p{\linewidth}}}]{
                \texttt{\tiny{[GM$|$GM]}}
                \texttt{Pénalité : (2,3) n'est pas vide, mais contient 'C'.} \\
\\ 
\texttt{Vous avez collectivement accumulé 2 pénalités sur 8.} \\
\texttt{Veuillez réessayer !} \\
            }
        }
    }
    & & \\ \\

    \theutterance \stepcounter{utterance}  
    & & \multicolumn{4}{p{0.6\linewidth}}{
        \cellcolor[rgb]{0.9,0.9,0.9}{
            \makecell[{{p{\linewidth}}}]{
                \texttt{\tiny{[GM$\rangle$P2]}}
                \texttt{Pénalité : (2,3) n'est pas vide, mais contient 'C'.} \\
\\ 
\texttt{Vous avez collectivement accumulé 2 pénalités sur 8.} \\
\texttt{Veuillez réessayer !} \\
            }
        }
    }
    & \\ \\

    \theutterance \stepcounter{utterance}  
    & & & \multicolumn{4}{p{0.6\linewidth}}{
        \cellcolor[rgb]{0.9,0.9,0.9}{
            \makecell[{{p{\linewidth}}}]{
                \texttt{\tiny{[GM$\langle$P2]}}
                \texttt{DIRE: J'ai reçu vos coordonnées mais je vois toujours un 'C' à (2,3). Pourriez{-}vous vérifier à nouveau la position exacte de votre 'C' et m'indiquer si c'est bien (2,3) ou une autre position ?} \\
            }
        }
    }
     \\ \\

    \theutterance \stepcounter{utterance}  
    & & & \multicolumn{2}{p{0.3\linewidth}}{
        \cellcolor[rgb]{0.9,0.9,0.9}{
            \makecell[{{p{\linewidth}}}]{
                \texttt{\tiny{[GM$|$GM]}}
                \texttt{DIRE: J'ai reçu vos coordonnées mais je vois toujours un 'C' à (2,3). Pourriez{-}vous vérifier à nouveau la position exacte de votre 'C' et m'indiquer si c'est bien (2,3) ou une autre position ?} \\
            }
        }
    }
    & & \\ \\

    \theutterance \stepcounter{utterance}  
    & & \multicolumn{4}{p{0.6\linewidth}}{
        \cellcolor[rgb]{0.9,0.9,0.9}{
            \makecell[{{p{\linewidth}}}]{
                \texttt{\tiny{[P1$\langle$GM]}}
                \texttt{Votre message a été transmis à l'autre joueur.} \\
\\ 
\\ 
\texttt{Vous jouez actuellement le tour 4 sur un maximum de 12.} \\
\texttt{Vous avez collectivement accumulé 2 pénalités sur 8.} \\
\texttt{L'autre joueur a envoyé ce message :} \\
\\ 
\\ 
\texttt{"J'ai reçu vos coordonnées mais je vois toujours un 'C' à (2,3). Pourriez{-}vous vérifier à nouveau la position exacte de votre 'C' et m'indiquer si c'est bien (2,3) ou une autre position ?"} \\
\\ 
\\ 
\texttt{Quelle est votre prochaine commande ?} \\
            }
        }
    }
    & \\ \\

    \theutterance \stepcounter{utterance}  
    & \multicolumn{4}{p{0.6\linewidth}}{
        \cellcolor[rgb]{0.9,0.9,0.9}{
            \makecell[{{p{\linewidth}}}]{
                \texttt{\tiny{[P1$\rangle$GM]}}
                \texttt{DIRE: Pourriez{-}vous vérifier votre grille à nouveau pour la position exacte de votre 'C' ? Est{-}ce correctement à (2, 3) ou à une autre position ? Cela me permettra de planifier correctement mon propre déplacement du 'C'.} \\
            }
        }
    }
    & & \\ \\

    \theutterance \stepcounter{utterance}  
    & & & \multicolumn{2}{p{0.3\linewidth}}{
        \cellcolor[rgb]{0.9,0.9,0.9}{
            \makecell[{{p{\linewidth}}}]{
                \texttt{\tiny{[GM$|$GM]}}
                \texttt{DIRE: Pourriez{-}vous vérifier votre grille à nouveau pour la position exacte de votre 'C' ? Est{-}ce correctement à (2, 3) ou à une autre position ? Cela me permettra de planifier correctement mon propre déplacement du 'C'.} \\
            }
        }
    }
    & & \\ \\

    \theutterance \stepcounter{utterance}  
    & & \multicolumn{4}{p{0.6\linewidth}}{
        \cellcolor[rgb]{0.9,0.9,0.9}{
            \makecell[{{p{\linewidth}}}]{
                \texttt{\tiny{[GM$\rangle$P2]}}
                \texttt{Votre message a été transmis à l'autre joueur.} \\
\\ 
\\ 
\texttt{Vous jouez actuellement le tour 5 sur un maximum de 12.} \\
\texttt{Vous avez collectivement accumulé 2 pénalités sur 8.} \\
\texttt{L'autre joueur a envoyé ce message :} \\
\\ 
\\ 
\texttt{"Pourriez{-}vous vérifier votre grille à nouveau pour la position exacte de votre 'C' ? Est{-}ce correctement à (2, 3) ou à une autre position ? Cela me permettra de planifier correctement mon propre déplacement du 'C'."} \\
\\ 
\\ 
\texttt{Quelle est votre prochaine commande ?} \\
            }
        }
    }
    & \\ \\

    \theutterance \stepcounter{utterance}  
    & & & \multicolumn{4}{p{0.6\linewidth}}{
        \cellcolor[rgb]{0.9,0.9,0.9}{
            \makecell[{{p{\linewidth}}}]{
                \texttt{\tiny{[GM$\langle$P2]}}
                \texttt{DIRE: Sur ma grille, le 'C' se trouve actuellement à la position (2, 3). Je vais donc déplacer mon 'C' vers une autre position pour laisser place à celui de l'autre joueur. Je suggère de déplacer mon 'C' vers (3, 3).} \\
            }
        }
    }
     \\ \\

    \theutterance \stepcounter{utterance}  
    & & & \multicolumn{2}{p{0.3\linewidth}}{
        \cellcolor[rgb]{0.9,0.9,0.9}{
            \makecell[{{p{\linewidth}}}]{
                \texttt{\tiny{[GM$|$GM]}}
                \texttt{DIRE: Sur ma grille, le 'C' se trouve actuellement à la position (2, 3). Je vais donc déplacer mon 'C' vers une autre position pour laisser place à celui de l'autre joueur. Je suggère de déplacer mon 'C' vers (3, 3).} \\
            }
        }
    }
    & & \\ \\

    \theutterance \stepcounter{utterance}  
    & & \multicolumn{4}{p{0.6\linewidth}}{
        \cellcolor[rgb]{0.9,0.9,0.9}{
            \makecell[{{p{\linewidth}}}]{
                \texttt{\tiny{[P1$\langle$GM]}}
                \texttt{Votre message a été transmis à l'autre joueur.} \\
\\ 
\\ 
\texttt{Vous jouez actuellement le tour 5 sur un maximum de 12.} \\
\texttt{Vous avez collectivement accumulé 2 pénalités sur 8.} \\
\texttt{L'autre joueur a envoyé ce message :} \\
\\ 
\\ 
\texttt{"Sur ma grille, le 'C' se trouve actuellement à la position (2, 3). Je vais donc déplacer mon 'C' vers une autre position pour laisser place à celui de l'autre joueur. Je suggère de déplacer mon 'C' vers (3, 3)."} \\
\\ 
\\ 
\texttt{Quelle est votre prochaine commande ?} \\
            }
        }
    }
    & \\ \\

    \theutterance \stepcounter{utterance}  
    & \multicolumn{4}{p{0.6\linewidth}}{
        \cellcolor[rgb]{0.9,0.9,0.9}{
            \makecell[{{p{\linewidth}}}]{
                \texttt{\tiny{[P1$\rangle$GM]}}
                \texttt{DÉPLACER: C, (3, 3)} \\
            }
        }
    }
    & & \\ \\

    \theutterance \stepcounter{utterance}  
    & & & \multicolumn{2}{p{0.3\linewidth}}{
        \cellcolor[rgb]{0.9,0.9,0.9}{
            \makecell[{{p{\linewidth}}}]{
                \texttt{\tiny{[GM$|$GM]}}
                \texttt{DÉPLACER: C, (3, 3)} \\
            }
        }
    }
    & & \\ \\

    \theutterance \stepcounter{utterance}  
    & & & \multicolumn{2}{p{0.3\linewidth}}{
        \cellcolor[rgb]{0.9,0.9,0.9}{
            \makecell[{{p{\linewidth}}}]{
                \texttt{\tiny{[GM$|$GM]}}
                \texttt{C moved to (3, 3)} \\
            }
        }
    }
    & & \\ \\

    \theutterance \stepcounter{utterance}  
    & & \multicolumn{4}{p{0.6\linewidth}}{
        \cellcolor[rgb]{0.9,0.9,0.9}{
            \makecell[{{p{\linewidth}}}]{
                \texttt{\tiny{[GM$\rangle$P2]}}
                \texttt{Votre message a été transmis à l'autre joueur.} \\
\\ 
\\ 
\texttt{Vous jouez actuellement le tour 6 sur un maximum de 12.} \\
\texttt{Vous avez collectivement accumulé 2 pénalités sur 8.} \\
\texttt{L'autre joueur a déplacé un objet sur sa grille.} \\
\\ 
\\ 
\texttt{Quelle est votre prochaine commande ?} \\
            }
        }
    }
    & \\ \\

    \theutterance \stepcounter{utterance}  
    & & & \multicolumn{4}{p{0.6\linewidth}}{
        \cellcolor[rgb]{0.9,0.9,0.9}{
            \makecell[{{p{\linewidth}}}]{
                \texttt{\tiny{[GM$\langle$P2]}}
                \texttt{DIRE: Je constate que vous avez déplacé votre 'C' vers (2,3). Mon 'C' est maintenant déplacé vers (3,3) pour libérer la position (2,3). Pourriez{-}vous vérifier si, sur votre grille, le 'C' est bien placé à (2,3) ou si nous devons continuer à harmoniser les positions ?} \\
            }
        }
    }
     \\ \\

    \theutterance \stepcounter{utterance}  
    & & & \multicolumn{2}{p{0.3\linewidth}}{
        \cellcolor[rgb]{0.9,0.9,0.9}{
            \makecell[{{p{\linewidth}}}]{
                \texttt{\tiny{[GM$|$GM]}}
                \texttt{DIRE: Je constate que vous avez déplacé votre 'C' vers (2,3). Mon 'C' est maintenant déplacé vers (3,3) pour libérer la position (2,3). Pourriez{-}vous vérifier si, sur votre grille, le 'C' est bien placé à (2,3) ou si nous devons continuer à harmoniser les positions ?} \\
            }
        }
    }
    & & \\ \\

    \theutterance \stepcounter{utterance}  
    & & \multicolumn{4}{p{0.6\linewidth}}{
        \cellcolor[rgb]{0.9,0.9,0.9}{
            \makecell[{{p{\linewidth}}}]{
                \texttt{\tiny{[P1$\langle$GM]}}
                \texttt{Déplacement de 'C' vers (3,3) réussi. Votre grille mise à jour ressemble à ceci :} \\
\texttt{```} \\
\texttt{ 1234567} \\
\texttt{╔═══════╗} \\
\texttt{║◌◌◌◌◌◌◌║ 1} \\
\texttt{║◌◌◌◌◌◌◌║ 2} \\
\texttt{║◌◌C◌◌┌─╢ 3} \\
\texttt{║◌◌◌┌─┤◌║ 4} \\
\texttt{║◌◌◌│L├─╢ 5} \\
\texttt{╟───┼─┘◌║ 6} \\
\texttt{║◌◌◌│◌◌P║ 7} \\
\texttt{╚═══╧═══╝} \\
\texttt{```} \\
\texttt{Vous jouez actuellement le tour 6 sur un maximum de 12.} \\
\texttt{Vous avez collectivement accumulé 2 pénalités sur 8.} \\
\texttt{L'autre joueur a envoyé ce message :} \\
\\ 
\\ 
\texttt{"Je constate que vous avez déplacé votre 'C' vers (2,3). Mon 'C' est maintenant déplacé vers (3,3) pour libérer la position (2,3). Pourriez{-}vous vérifier si, sur votre grille, le 'C' est bien placé à (2,3) ou si nous devons continuer à harmoniser les positions ?"} \\
\\ 
\\ 
\texttt{Quelle est votre prochaine commande ?} \\
            }
        }
    }
    & \\ \\

    \theutterance \stepcounter{utterance}  
    & \multicolumn{4}{p{0.6\linewidth}}{
        \cellcolor[rgb]{0.9,0.9,0.9}{
            \makecell[{{p{\linewidth}}}]{
                \texttt{\tiny{[P1$\rangle$GM]}}
                \texttt{DIRE: J'ai bien déplacé mon 'C' vers (2,3). Pourriez{-}vous déplacer votre 'C' vers (3,3) maintenant ? Je vais vérifier si cette position est libre sur votre grille en confirmant votre déplacement.} \\
            }
        }
    }
    & & \\ \\

    \theutterance \stepcounter{utterance}  
    & & & \multicolumn{2}{p{0.3\linewidth}}{
        \cellcolor[rgb]{0.9,0.9,0.9}{
            \makecell[{{p{\linewidth}}}]{
                \texttt{\tiny{[GM$|$GM]}}
                \texttt{DIRE: J'ai bien déplacé mon 'C' vers (2,3). Pourriez{-}vous déplacer votre 'C' vers (3,3) maintenant ? Je vais vérifier si cette position est libre sur votre grille en confirmant votre déplacement.} \\
            }
        }
    }
    & & \\ \\

    \theutterance \stepcounter{utterance}  
    & & \multicolumn{4}{p{0.6\linewidth}}{
        \cellcolor[rgb]{0.9,0.9,0.9}{
            \makecell[{{p{\linewidth}}}]{
                \texttt{\tiny{[GM$\rangle$P2]}}
                \texttt{Votre message a été transmis à l'autre joueur.} \\
\\ 
\\ 
\texttt{Vous jouez actuellement le tour 7 sur un maximum de 12.} \\
\texttt{Vous avez collectivement accumulé 2 pénalités sur 8.} \\
\texttt{L'autre joueur a envoyé ce message :} \\
\\ 
\\ 
\texttt{"J'ai bien déplacé mon 'C' vers (2,3). Pourriez{-}vous déplacer votre 'C' vers (3,3) maintenant ? Je vais vérifier si cette position est libre sur votre grille en confirmant votre déplacement."} \\
\\ 
\\ 
\texttt{Quelle est votre prochaine commande ?} \\
            }
        }
    }
    & \\ \\

    \theutterance \stepcounter{utterance}  
    & & & \multicolumn{4}{p{0.6\linewidth}}{
        \cellcolor[rgb]{0.9,0.9,0.9}{
            \makecell[{{p{\linewidth}}}]{
                \texttt{\tiny{[GM$\langle$P2]}}
                \texttt{DÉPLACER: C, (3, 3)} \\
            }
        }
    }
     \\ \\

    \theutterance \stepcounter{utterance}  
    & & & \multicolumn{2}{p{0.3\linewidth}}{
        \cellcolor[rgb]{0.9,0.9,0.9}{
            \makecell[{{p{\linewidth}}}]{
                \texttt{\tiny{[GM$|$GM]}}
                \texttt{DÉPLACER: C, (3, 3)} \\
            }
        }
    }
    & & \\ \\

    \theutterance \stepcounter{utterance}  
    & & & \multicolumn{2}{p{0.3\linewidth}}{
        \cellcolor[rgb]{0.9,0.9,0.9}{
            \makecell[{{p{\linewidth}}}]{
                \texttt{\tiny{[GM$|$GM]}}
                \texttt{C moved to (3, 3)} \\
            }
        }
    }
    & & \\ \\

    \theutterance \stepcounter{utterance}  
    & & \multicolumn{4}{p{0.6\linewidth}}{
        \cellcolor[rgb]{0.9,0.9,0.9}{
            \makecell[{{p{\linewidth}}}]{
                \texttt{\tiny{[P1$\langle$GM]}}
                \texttt{Votre message a été transmis à l'autre joueur.} \\
\\ 
\\ 
\texttt{Vous jouez actuellement le tour 7 sur un maximum de 12.} \\
\texttt{Vous avez collectivement accumulé 2 pénalités sur 8.} \\
\texttt{L'autre joueur a déplacé un objet sur sa grille.} \\
\\ 
\\ 
\texttt{Quelle est votre prochaine commande ?} \\
            }
        }
    }
    & \\ \\

    \theutterance \stepcounter{utterance}  
    & \multicolumn{4}{p{0.6\linewidth}}{
        \cellcolor[rgb]{0.9,0.9,0.9}{
            \makecell[{{p{\linewidth}}}]{
                \texttt{\tiny{[P1$\rangle$GM]}}
                \texttt{DIRE: Avez{-}vous pu déplacer votre 'C' vers (3,3) ? Si oui, pourriez{-}vous me le confirmer ? Sinon, pourriez{-}vous me donner les nouvelles coordonnées de votre 'C' après le déplacement pour que je puisse ajuster ma stratégie en conséquence ?} \\
            }
        }
    }
    & & \\ \\

    \theutterance \stepcounter{utterance}  
    & & & \multicolumn{2}{p{0.3\linewidth}}{
        \cellcolor[rgb]{0.9,0.9,0.9}{
            \makecell[{{p{\linewidth}}}]{
                \texttt{\tiny{[GM$|$GM]}}
                \texttt{DIRE: Avez{-}vous pu déplacer votre 'C' vers (3,3) ? Si oui, pourriez{-}vous me le confirmer ? Sinon, pourriez{-}vous me donner les nouvelles coordonnées de votre 'C' après le déplacement pour que je puisse ajuster ma stratégie en conséquence ?} \\
            }
        }
    }
    & & \\ \\

    \theutterance \stepcounter{utterance}  
    & & \multicolumn{4}{p{0.6\linewidth}}{
        \cellcolor[rgb]{0.9,0.9,0.9}{
            \makecell[{{p{\linewidth}}}]{
                \texttt{\tiny{[GM$\rangle$P2]}}
                \texttt{Déplacement de 'C' vers (3,3) réussi. Votre grille mise à jour ressemble à ceci :} \\
\texttt{```} \\
\texttt{ 1234567} \\
\texttt{╔═══════╗} \\
\texttt{║◌◌◌◌◌◌◌║ 1} \\
\texttt{║◌◌◌◌◌◌◌║ 2} \\
\texttt{║◌◌C◌◌┌─╢ 3} \\
\texttt{║◌◌◌┌─┤◌║ 4} \\
\texttt{║◌P◌│◌├─╢ 5} \\
\texttt{╟───┼─┘◌║ 6} \\
\texttt{║L◌◌│◌◌◌║ 7} \\
\texttt{╚═══╧═══╝} \\
\texttt{```} \\
\texttt{Vous jouez actuellement le tour 8 sur un maximum de 12.} \\
\texttt{Vous avez collectivement accumulé 2 pénalités sur 8.} \\
\texttt{L'autre joueur a envoyé ce message :} \\
\\ 
\\ 
\texttt{"Avez{-}vous pu déplacer votre 'C' vers (3,3) ? Si oui, pourriez{-}vous me le confirmer ? Sinon, pourriez{-}vous me donner les nouvelles coordonnées de votre 'C' après le déplacement pour que je puisse ajuster ma stratégie en conséquence ?"} \\
\\ 
\\ 
\texttt{Quelle est votre prochaine commande ?} \\
            }
        }
    }
    & \\ \\

    \theutterance \stepcounter{utterance}  
    & & & \multicolumn{4}{p{0.6\linewidth}}{
        \cellcolor[rgb]{0.9,0.9,0.9}{
            \makecell[{{p{\linewidth}}}]{
                \texttt{\tiny{[GM$\langle$P2]}}
                \texttt{DIRE: Mon 'C' a bien été déplacé vers (3,3). Veuillez me confirmer que votre 'C' est bien placé à (2,3) pour que nous puissions continuer à harmoniser nos configurations.} \\
            }
        }
    }
     \\ \\

    \theutterance \stepcounter{utterance}  
    & & & \multicolumn{2}{p{0.3\linewidth}}{
        \cellcolor[rgb]{0.9,0.9,0.9}{
            \makecell[{{p{\linewidth}}}]{
                \texttt{\tiny{[GM$|$GM]}}
                \texttt{DIRE: Mon 'C' a bien été déplacé vers (3,3). Veuillez me confirmer que votre 'C' est bien placé à (2,3) pour que nous puissions continuer à harmoniser nos configurations.} \\
            }
        }
    }
    & & \\ \\

    \theutterance \stepcounter{utterance}  
    & & \multicolumn{4}{p{0.6\linewidth}}{
        \cellcolor[rgb]{0.9,0.9,0.9}{
            \makecell[{{p{\linewidth}}}]{
                \texttt{\tiny{[P1$\langle$GM]}}
                \texttt{Votre message a été transmis à l'autre joueur.} \\
\\ 
\\ 
\texttt{Vous jouez actuellement le tour 8 sur un maximum de 12.} \\
\texttt{Vous avez collectivement accumulé 2 pénalités sur 8.} \\
\texttt{L'autre joueur a envoyé ce message :} \\
\\ 
\\ 
\texttt{"Mon 'C' a bien été déplacé vers (3,3). Veuillez me confirmer que votre 'C' est bien placé à (2,3) pour que nous puissions continuer à harmoniser nos configurations."} \\
\\ 
\\ 
\texttt{Quelle est votre prochaine commande ?} \\
            }
        }
    }
    & \\ \\

    \theutterance \stepcounter{utterance}  
    & \multicolumn{4}{p{0.6\linewidth}}{
        \cellcolor[rgb]{0.9,0.9,0.9}{
            \makecell[{{p{\linewidth}}}]{
                \texttt{\tiny{[P1$\rangle$GM]}}
                \texttt{DIRE: Confirmé, mon 'C' est bien placé à (2,3). Pour harmoniser nos configurations, je suggère de déplacer mon 'L' vers (5, 4). Est{-}ce que cette position est libre sur votre grille ?} \\
            }
        }
    }
    & & \\ \\

    \theutterance \stepcounter{utterance}  
    & & & \multicolumn{2}{p{0.3\linewidth}}{
        \cellcolor[rgb]{0.9,0.9,0.9}{
            \makecell[{{p{\linewidth}}}]{
                \texttt{\tiny{[GM$|$GM]}}
                \texttt{DIRE: Confirmé, mon 'C' est bien placé à (2,3). Pour harmoniser nos configurations, je suggère de déplacer mon 'L' vers (5, 4). Est{-}ce que cette position est libre sur votre grille ?} \\
            }
        }
    }
    & & \\ \\

    \theutterance \stepcounter{utterance}  
    & & \multicolumn{4}{p{0.6\linewidth}}{
        \cellcolor[rgb]{0.9,0.9,0.9}{
            \makecell[{{p{\linewidth}}}]{
                \texttt{\tiny{[GM$\rangle$P2]}}
                \texttt{Votre message a été transmis à l'autre joueur.} \\
\\ 
\\ 
\texttt{Vous jouez actuellement le tour 9 sur un maximum de 12.} \\
\texttt{Vous avez collectivement accumulé 2 pénalités sur 8.} \\
\texttt{L'autre joueur a envoyé ce message :} \\
\\ 
\\ 
\texttt{"Confirmé, mon 'C' est bien placé à (2,3). Pour harmoniser nos configurations, je suggère de déplacer mon 'L' vers (5, 4). Est{-}ce que cette position est libre sur votre grille ?"} \\
\\ 
\\ 
\texttt{Quelle est votre prochaine commande ?} \\
            }
        }
    }
    & \\ \\

    \theutterance \stepcounter{utterance}  
    & & & \multicolumn{4}{p{0.6\linewidth}}{
        \cellcolor[rgb]{0.9,0.9,0.9}{
            \makecell[{{p{\linewidth}}}]{
                \texttt{\tiny{[GM$\langle$P2]}}
                \texttt{DIRE: Je vérifie la position (5,4) sur ma grille. Elle est actuellement libre et convient pour le déplacement de votre 'L'. Veuillez confirmer si cette position est adaptée ou si vous préférez une autre position.} \\
            }
        }
    }
     \\ \\

    \theutterance \stepcounter{utterance}  
    & & & \multicolumn{2}{p{0.3\linewidth}}{
        \cellcolor[rgb]{0.9,0.9,0.9}{
            \makecell[{{p{\linewidth}}}]{
                \texttt{\tiny{[GM$|$GM]}}
                \texttt{DIRE: Je vérifie la position (5,4) sur ma grille. Elle est actuellement libre et convient pour le déplacement de votre 'L'. Veuillez confirmer si cette position est adaptée ou si vous préférez une autre position.} \\
            }
        }
    }
    & & \\ \\

    \theutterance \stepcounter{utterance}  
    & & \multicolumn{4}{p{0.6\linewidth}}{
        \cellcolor[rgb]{0.9,0.9,0.9}{
            \makecell[{{p{\linewidth}}}]{
                \texttt{\tiny{[P1$\langle$GM]}}
                \texttt{Votre message a été transmis à l'autre joueur.} \\
\\ 
\\ 
\texttt{Vous jouez actuellement le tour 9 sur un maximum de 12.} \\
\texttt{Vous avez collectivement accumulé 2 pénalités sur 8.} \\
\texttt{L'autre joueur a envoyé ce message :} \\
\\ 
\\ 
\texttt{"Je vérifie la position (5,4) sur ma grille. Elle est actuellement libre et convient pour le déplacement de votre 'L'. Veuillez confirmer si cette position est adaptée ou si vous préférez une autre position."} \\
\\ 
\\ 
\texttt{Quelle est votre prochaine commande ?} \\
            }
        }
    }
    & \\ \\

    \theutterance \stepcounter{utterance}  
    & \multicolumn{4}{p{0.6\linewidth}}{
        \cellcolor[rgb]{0.9,0.9,0.9}{
            \makecell[{{p{\linewidth}}}]{
                \texttt{\tiny{[P1$\rangle$GM]}}
                \texttt{DÉPLACER: L, (5,4)} \\
            }
        }
    }
    & & \\ \\

    \theutterance \stepcounter{utterance}  
    & & & \multicolumn{2}{p{0.3\linewidth}}{
        \cellcolor[rgb]{0.9,0.9,0.9}{
            \makecell[{{p{\linewidth}}}]{
                \texttt{\tiny{[GM$|$GM]}}
                \texttt{DÉPLACER: L, (5,4)} \\
            }
        }
    }
    & & \\ \\

    \theutterance \stepcounter{utterance}  
    & & & \multicolumn{2}{p{0.3\linewidth}}{
        \cellcolor[rgb]{0.9,0.9,0.9}{
            \makecell[{{p{\linewidth}}}]{
                \texttt{\tiny{[GM$|$GM]}}
                \texttt{Pénalité : (5,4) n'est pas vide, mais contient '─'.} \\
\\ 
\texttt{Vous avez collectivement accumulé 3 pénalités sur 8.} \\
\texttt{Veuillez réessayer !} \\
            }
        }
    }
    & & \\ \\

    \theutterance \stepcounter{utterance}  
    & & \multicolumn{4}{p{0.6\linewidth}}{
        \cellcolor[rgb]{0.9,0.9,0.9}{
            \makecell[{{p{\linewidth}}}]{
                \texttt{\tiny{[P1$\langle$GM]}}
                \texttt{Pénalité : (5,4) n'est pas vide, mais contient '─'.} \\
\\ 
\texttt{Vous avez collectivement accumulé 3 pénalités sur 8.} \\
\texttt{Veuillez réessayer !} \\
            }
        }
    }
    & \\ \\

    \theutterance \stepcounter{utterance}  
    & \multicolumn{4}{p{0.6\linewidth}}{
        \cellcolor[rgb]{0.9,0.9,0.9}{
            \makecell[{{p{\linewidth}}}]{
                \texttt{\tiny{[P1$\rangle$GM]}}
                \texttt{DIRE: Je vois que (5,4) contient un '─' sur votre grille, je m'excuse pour la confusion. Pourriez{-}vous vérifier les coordonnées exactes où vous pourriez placer mon 'L' pour qu'il soit libre ? Essayez (5, 5) à la place.} \\
            }
        }
    }
    & & \\ \\

    \theutterance \stepcounter{utterance}  
    & & & \multicolumn{2}{p{0.3\linewidth}}{
        \cellcolor[rgb]{0.9,0.9,0.9}{
            \makecell[{{p{\linewidth}}}]{
                \texttt{\tiny{[GM$|$GM]}}
                \texttt{DIRE: Je vois que (5,4) contient un '─' sur votre grille, je m'excuse pour la confusion. Pourriez{-}vous vérifier les coordonnées exactes où vous pourriez placer mon 'L' pour qu'il soit libre ? Essayez (5, 5) à la place.} \\
            }
        }
    }
    & & \\ \\

    \theutterance \stepcounter{utterance}  
    & & \multicolumn{4}{p{0.6\linewidth}}{
        \cellcolor[rgb]{0.9,0.9,0.9}{
            \makecell[{{p{\linewidth}}}]{
                \texttt{\tiny{[GM$\rangle$P2]}}
                \texttt{Votre message a été transmis à l'autre joueur.} \\
\\ 
\\ 
\texttt{Vous jouez actuellement le tour 10 sur un maximum de 12.} \\
\texttt{Vous avez collectivement accumulé 3 pénalités sur 8.} \\
\texttt{L'autre joueur a envoyé ce message :} \\
\\ 
\\ 
\texttt{"Je vois que (5,4) contient un '─' sur votre grille, je m'excuse pour la confusion. Pourriez{-}vous vérifier les coordonnées exactes où vous pourriez placer mon 'L' pour qu'il soit libre ? Essayez (5, 5) à la place."} \\
\\ 
\\ 
\texttt{Quelle est votre prochaine commande ?} \\
            }
        }
    }
    & \\ \\

    \theutterance \stepcounter{utterance}  
    & & & \multicolumn{4}{p{0.6\linewidth}}{
        \cellcolor[rgb]{0.9,0.9,0.9}{
            \makecell[{{p{\linewidth}}}]{
                \texttt{\tiny{[GM$\langle$P2]}}
                \texttt{DIRE: Sur ma grille, la position (5,5) semble être libre. Je vais déplacer votre 'L' vers (5,5) pour harmoniser nos configurations.} \\
            }
        }
    }
     \\ \\

    \theutterance \stepcounter{utterance}  
    & & & \multicolumn{2}{p{0.3\linewidth}}{
        \cellcolor[rgb]{0.9,0.9,0.9}{
            \makecell[{{p{\linewidth}}}]{
                \texttt{\tiny{[GM$|$GM]}}
                \texttt{DIRE: Sur ma grille, la position (5,5) semble être libre. Je vais déplacer votre 'L' vers (5,5) pour harmoniser nos configurations.} \\
            }
        }
    }
    & & \\ \\

    \theutterance \stepcounter{utterance}  
    & & \multicolumn{4}{p{0.6\linewidth}}{
        \cellcolor[rgb]{0.9,0.9,0.9}{
            \makecell[{{p{\linewidth}}}]{
                \texttt{\tiny{[P1$\langle$GM]}}
                \texttt{Votre message a été transmis à l'autre joueur.} \\
\\ 
\\ 
\texttt{Vous jouez actuellement le tour 10 sur un maximum de 12.} \\
\texttt{Vous avez collectivement accumulé 3 pénalités sur 8.} \\
\texttt{L'autre joueur a envoyé ce message :} \\
\\ 
\\ 
\texttt{"Sur ma grille, la position (5,5) semble être libre. Je vais déplacer votre 'L' vers (5,5) pour harmoniser nos configurations."} \\
\\ 
\\ 
\texttt{Quelle est votre prochaine commande ?} \\
            }
        }
    }
    & \\ \\

    \theutterance \stepcounter{utterance}  
    & \multicolumn{4}{p{0.6\linewidth}}{
        \cellcolor[rgb]{0.9,0.9,0.9}{
            \makecell[{{p{\linewidth}}}]{
                \texttt{\tiny{[P1$\rangle$GM]}}
                \texttt{DÉPLACER: L, (5,5)} \\
            }
        }
    }
    & & \\ \\

    \theutterance \stepcounter{utterance}  
    & & & \multicolumn{2}{p{0.3\linewidth}}{
        \cellcolor[rgb]{0.9,0.9,0.9}{
            \makecell[{{p{\linewidth}}}]{
                \texttt{\tiny{[GM$|$GM]}}
                \texttt{DÉPLACER: L, (5,5)} \\
            }
        }
    }
    & & \\ \\

    \theutterance \stepcounter{utterance}  
    & & & \multicolumn{2}{p{0.3\linewidth}}{
        \cellcolor[rgb]{0.9,0.9,0.9}{
            \makecell[{{p{\linewidth}}}]{
                \texttt{\tiny{[GM$|$GM]}}
                \texttt{Pénalité : (5,5) n'est pas vide, mais contient 'L'.} \\
\\ 
\texttt{Vous avez collectivement accumulé 4 pénalités sur 8.} \\
\texttt{Veuillez réessayer !} \\
            }
        }
    }
    & & \\ \\

    \theutterance \stepcounter{utterance}  
    & & \multicolumn{4}{p{0.6\linewidth}}{
        \cellcolor[rgb]{0.9,0.9,0.9}{
            \makecell[{{p{\linewidth}}}]{
                \texttt{\tiny{[P1$\langle$GM]}}
                \texttt{Pénalité : (5,5) n'est pas vide, mais contient 'L'.} \\
\\ 
\texttt{Vous avez collectivement accumulé 4 pénalités sur 8.} \\
\texttt{Veuillez réessayer !} \\
            }
        }
    }
    & \\ \\

    \theutterance \stepcounter{utterance}  
    & \multicolumn{4}{p{0.6\linewidth}}{
        \cellcolor[rgb]{0.9,0.9,0.9}{
            \makecell[{{p{\linewidth}}}]{
                \texttt{\tiny{[P1$\rangle$GM]}}
                \texttt{DIRE: Je suis désolé, (5,5) contient un 'L'. Pourriez{-}vous vérifier les coordonnées où vous pourriez déplacer ma 'L' pour qu'elle soit libre ? Essayez (5, 6) ou (4, 5).} \\
            }
        }
    }
    & & \\ \\

    \theutterance \stepcounter{utterance}  
    & & & \multicolumn{2}{p{0.3\linewidth}}{
        \cellcolor[rgb]{0.9,0.9,0.9}{
            \makecell[{{p{\linewidth}}}]{
                \texttt{\tiny{[GM$|$GM]}}
                \texttt{DIRE: Je suis désolé, (5,5) contient un 'L'. Pourriez{-}vous vérifier les coordonnées où vous pourriez déplacer ma 'L' pour qu'elle soit libre ? Essayez (5, 6) ou (4, 5).} \\
            }
        }
    }
    & & \\ \\

    \theutterance \stepcounter{utterance}  
    & & \multicolumn{4}{p{0.6\linewidth}}{
        \cellcolor[rgb]{0.9,0.9,0.9}{
            \makecell[{{p{\linewidth}}}]{
                \texttt{\tiny{[GM$\rangle$P2]}}
                \texttt{Votre message a été transmis à l'autre joueur.} \\
\\ 
\\ 
\texttt{Vous jouez actuellement le tour 11 sur un maximum de 12.} \\
\texttt{Vous avez collectivement accumulé 4 pénalités sur 8.} \\
\texttt{L'autre joueur a envoyé ce message :} \\
\\ 
\\ 
\texttt{"Je suis désolé, (5,5) contient un 'L'. Pourriez{-}vous vérifier les coordonnées où vous pourriez déplacer ma 'L' pour qu'elle soit libre ? Essayez (5, 6) ou (4, 5)."} \\
\\ 
\\ 
\texttt{Quelle est votre prochaine commande ?} \\
            }
        }
    }
    & \\ \\

    \theutterance \stepcounter{utterance}  
    & & & \multicolumn{4}{p{0.6\linewidth}}{
        \cellcolor[rgb]{0.9,0.9,0.9}{
            \makecell[{{p{\linewidth}}}]{
                \texttt{\tiny{[GM$\langle$P2]}}
                \texttt{DIRE: Sur ma grille, la position (4, 5) semble être libre. Je vais déplacer votre 'L' vers (4, 5) pour harmoniser les configurations des deux plateaux de jeu.} \\
            }
        }
    }
     \\ \\

    \theutterance \stepcounter{utterance}  
    & & & \multicolumn{2}{p{0.3\linewidth}}{
        \cellcolor[rgb]{0.9,0.9,0.9}{
            \makecell[{{p{\linewidth}}}]{
                \texttt{\tiny{[GM$|$GM]}}
                \texttt{DIRE: Sur ma grille, la position (4, 5) semble être libre. Je vais déplacer votre 'L' vers (4, 5) pour harmoniser les configurations des deux plateaux de jeu.} \\
            }
        }
    }
    & & \\ \\

    \theutterance \stepcounter{utterance}  
    & & \multicolumn{4}{p{0.6\linewidth}}{
        \cellcolor[rgb]{0.9,0.9,0.9}{
            \makecell[{{p{\linewidth}}}]{
                \texttt{\tiny{[P1$\langle$GM]}}
                \texttt{Votre message a été transmis à l'autre joueur.} \\
\\ 
\\ 
\texttt{Vous jouez actuellement le tour 11 sur un maximum de 12.} \\
\texttt{Vous avez collectivement accumulé 4 pénalités sur 8.} \\
\texttt{L'autre joueur a envoyé ce message :} \\
\\ 
\\ 
\texttt{"Sur ma grille, la position (4, 5) semble être libre. Je vais déplacer votre 'L' vers (4, 5) pour harmoniser les configurations des deux plateaux de jeu."} \\
\\ 
\\ 
\texttt{Quelle est votre prochaine commande ?} \\
            }
        }
    }
    & \\ \\

    \theutterance \stepcounter{utterance}  
    & \multicolumn{4}{p{0.6\linewidth}}{
        \cellcolor[rgb]{0.9,0.9,0.9}{
            \makecell[{{p{\linewidth}}}]{
                \texttt{\tiny{[P1$\rangle$GM]}}
                \texttt{DÉPLACER: L, (4,5)} \\
            }
        }
    }
    & & \\ \\

    \theutterance \stepcounter{utterance}  
    & & & \multicolumn{2}{p{0.3\linewidth}}{
        \cellcolor[rgb]{0.9,0.9,0.9}{
            \makecell[{{p{\linewidth}}}]{
                \texttt{\tiny{[GM$|$GM]}}
                \texttt{DÉPLACER: L, (4,5)} \\
            }
        }
    }
    & & \\ \\

    \theutterance \stepcounter{utterance}  
    & & & \multicolumn{2}{p{0.3\linewidth}}{
        \cellcolor[rgb]{0.9,0.9,0.9}{
            \makecell[{{p{\linewidth}}}]{
                \texttt{\tiny{[GM$|$GM]}}
                \texttt{Pénalité : (4,5) n'est pas vide, mais contient '│'.} \\
\\ 
\texttt{Vous avez collectivement accumulé 5 pénalités sur 8.} \\
\texttt{Veuillez réessayer !} \\
            }
        }
    }
    & & \\ \\

    \theutterance \stepcounter{utterance}  
    & & \multicolumn{4}{p{0.6\linewidth}}{
        \cellcolor[rgb]{0.9,0.9,0.9}{
            \makecell[{{p{\linewidth}}}]{
                \texttt{\tiny{[P1$\langle$GM]}}
                \texttt{Pénalité : (4,5) n'est pas vide, mais contient '│'.} \\
\\ 
\texttt{Vous avez collectivement accumulé 5 pénalités sur 8.} \\
\texttt{Veuillez réessayer !} \\
            }
        }
    }
    & \\ \\

    \theutterance \stepcounter{utterance}  
    & \multicolumn{4}{p{0.6\linewidth}}{
        \cellcolor[rgb]{0.9,0.9,0.9}{
            \makecell[{{p{\linewidth}}}]{
                \texttt{\tiny{[P1$\rangle$GM]}}
                \texttt{DIRE: Je vois que (4,5) contient '│' sur votre grille, ce qui signifie qu'elle n'est pas libre. Pourriez{-}vous vérifier la position exacte où vous pourriez déplacer ma 'L' sans collision ? Essayez (3, 5) ou (5, 4).} \\
            }
        }
    }
    & & \\ \\

    \theutterance \stepcounter{utterance}  
    & & & \multicolumn{2}{p{0.3\linewidth}}{
        \cellcolor[rgb]{0.9,0.9,0.9}{
            \makecell[{{p{\linewidth}}}]{
                \texttt{\tiny{[GM$|$GM]}}
                \texttt{DIRE: Je vois que (4,5) contient '│' sur votre grille, ce qui signifie qu'elle n'est pas libre. Pourriez{-}vous vérifier la position exacte où vous pourriez déplacer ma 'L' sans collision ? Essayez (3, 5) ou (5, 4).} \\
            }
        }
    }
    & & \\ \\

    \theutterance \stepcounter{utterance}  
    & & \multicolumn{4}{p{0.6\linewidth}}{
        \cellcolor[rgb]{0.9,0.9,0.9}{
            \makecell[{{p{\linewidth}}}]{
                \texttt{\tiny{[GM$\rangle$P2]}}
                \texttt{Votre message a été transmis à l'autre joueur.} \\
\\ 
\\ 
\texttt{Vous jouez actuellement le tour 12 sur un maximum de 12.} \\
\texttt{Vous avez collectivement accumulé 5 pénalités sur 8.} \\
\texttt{L'autre joueur a envoyé ce message :} \\
\\ 
\\ 
\texttt{"Je vois que (4,5) contient '│' sur votre grille, ce qui signifie qu'elle n'est pas libre. Pourriez{-}vous vérifier la position exacte où vous pourriez déplacer ma 'L' sans collision ? Essayez (3, 5) ou (5, 4)."} \\
\\ 
\\ 
\texttt{Quelle est votre prochaine commande ?} \\
            }
        }
    }
    & \\ \\

    \theutterance \stepcounter{utterance}  
    & & & \multicolumn{4}{p{0.6\linewidth}}{
        \cellcolor[rgb]{0.9,0.9,0.9}{
            \makecell[{{p{\linewidth}}}]{
                \texttt{\tiny{[GM$\langle$P2]}}
                \texttt{DIRE: Sur ma grille, la position (5,4) semble être libre. Je vais déplacer votre 'L' vers (5,4) pour harmoniser les configurations des deux plateaux de jeu.} \\
            }
        }
    }
     \\ \\

    \theutterance \stepcounter{utterance}  
    & & & \multicolumn{2}{p{0.3\linewidth}}{
        \cellcolor[rgb]{0.9,0.9,0.9}{
            \makecell[{{p{\linewidth}}}]{
                \texttt{\tiny{[GM$|$GM]}}
                \texttt{DIRE: Sur ma grille, la position (5,4) semble être libre. Je vais déplacer votre 'L' vers (5,4) pour harmoniser les configurations des deux plateaux de jeu.} \\
            }
        }
    }
    & & \\ \\

    \theutterance \stepcounter{utterance}  
    & & & \multicolumn{2}{p{0.3\linewidth}}{
        \cellcolor[rgb]{0.9,0.9,0.9}{
            \makecell[{{p{\linewidth}}}]{
                \texttt{\tiny{[GM$|$GM]}}
                \texttt{Maximum number of rounds reached} \\
            }
        }
    }
    & & \\ \\

    \theutterance \stepcounter{utterance}  
    & & & \multicolumn{2}{p{0.3\linewidth}}{
        \cellcolor[rgb]{0.9,0.9,0.9}{
            \makecell[{{p{\linewidth}}}]{
                \texttt{\tiny{[GM$|$GM]}}
                \texttt{Initial states:} \\
\texttt{```} \\
\texttt{Player 1:} \\
\texttt{ 1234567} \\
\texttt{╔═══════╗} \\
\texttt{║◌◌◌◌◌◌◌║ 1} \\
\texttt{║◌C◌◌◌◌◌║ 2} \\
\texttt{║◌◌◌◌◌┌─╢ 3} \\
\texttt{║◌◌◌┌─┤◌║ 4} \\
\texttt{║◌◌◌│L├─╢ 5} \\
\texttt{╟───┼─┘◌║ 6} \\
\texttt{║◌◌◌│◌◌P║ 7} \\
\texttt{╚═══╧═══╝} \\
\\ 
\\ 
\texttt{Player 2:} \\
\texttt{ 1234567} \\
\texttt{╔═══════╗} \\
\texttt{║◌◌◌◌◌◌◌║ 1} \\
\texttt{║◌◌C◌◌◌◌║ 2} \\
\texttt{║◌◌◌◌◌┌─╢ 3} \\
\texttt{║◌◌◌┌─┤◌║ 4} \\
\texttt{║◌P◌│◌├─╢ 5} \\
\texttt{╟───┼─┘◌║ 6} \\
\texttt{║L◌◌│◌◌◌║ 7} \\
\texttt{╚═══╧═══╝} \\
\\ 
\texttt{```} \\
            }
        }
    }
    & & \\ \\

    \theutterance \stepcounter{utterance}  
    & & & \multicolumn{2}{p{0.3\linewidth}}{
        \cellcolor[rgb]{0.9,0.9,0.9}{
            \makecell[{{p{\linewidth}}}]{
                \texttt{\tiny{[GM$|$GM]}}
                \texttt{End states:} \\
\texttt{```} \\
\texttt{Player 1:} \\
\texttt{ 1234567} \\
\texttt{╔═══════╗} \\
\texttt{║◌◌◌◌◌◌◌║ 1} \\
\texttt{║◌◌◌◌◌◌◌║ 2} \\
\texttt{║◌◌C◌◌┌─╢ 3} \\
\texttt{║◌◌◌┌─┤◌║ 4} \\
\texttt{║◌◌◌│L├─╢ 5} \\
\texttt{╟───┼─┘◌║ 6} \\
\texttt{║◌◌◌│◌◌P║ 7} \\
\texttt{╚═══╧═══╝} \\
\\ 
\\ 
\texttt{Player 2:} \\
\texttt{ 1234567} \\
\texttt{╔═══════╗} \\
\texttt{║◌◌◌◌◌◌◌║ 1} \\
\texttt{║◌◌◌◌◌◌◌║ 2} \\
\texttt{║◌◌C◌◌┌─╢ 3} \\
\texttt{║◌◌◌┌─┤◌║ 4} \\
\texttt{║◌P◌│◌├─╢ 5} \\
\texttt{╟───┼─┘◌║ 6} \\
\texttt{║L◌◌│◌◌◌║ 7} \\
\texttt{╚═══╧═══╝} \\
\\ 
\texttt{```} \\
            }
        }
    }
    & & \\ \\

    \theutterance \stepcounter{utterance}  
    & & & \multicolumn{2}{p{0.3\linewidth}}{
        \cellcolor[rgb]{0.9,0.9,0.9}{
            \makecell[{{p{\linewidth}}}]{
                \texttt{\tiny{[GM$|$GM]}}
                \texttt{* success: True} \\
\texttt{* lose: False} \\
\texttt{* aborted: False} \\
\texttt{{-}{-}{-}{-}{-}{-}{-}} \\
\texttt{* Shifts: 0.00} \\
\texttt{* Max Shifts: 4.00} \\
\texttt{* Min Shifts: 2.00} \\
\texttt{* End Distance Sum: 9.86} \\
\texttt{* Init Distance Sum: 10.86} \\
\texttt{* Expected Distance Sum: 12.57} \\
\texttt{* Penalties: 5.00} \\
\texttt{* Max Penalties: 8.00} \\
\texttt{* Rounds: 12.00} \\
\texttt{* Max Rounds: 12.00} \\
\texttt{* Object Count: 3.00} \\
            }
        }
    }
    & & \\ \\

\end{supertabular}
}

\end{document}
